% Introduction (PUB-REF-0001 AN-21 F-BOR: what the book covers, the manuals
% and their jobs with routing, installing, typefaces, contacting Initech).
% Sources (Law 1): ADR-0007 DEC-01, DEC-02, DEC-08; os/tps/tps.pas;
% docs/pubs/originals/borland-turbo-pascal.md s1, s3.
\Gates{test-pubs-build, test-compiler, test-pubs-tps}
\IntroductionChapter

This book is the language definition and the technical reference for Turbo
Initech\reg{}.  It tells you exactly which programs Turbo Initech accepts,
what each construct means, how the compiler lays out your data and your
routines in memory, what the run-time library does, and what every message
of the compiler means.  It is written for programmers who already know how
to write a program and want the precise rule.  If you are new to Pascal,
start with the \doctitle{User's Guide}, which teaches the language by
example.

\section{What Turbo Initech is}
Turbo Initech is a single-pass Pascal compiler for InitechOS.  It reads one
source file, checks it completely, and writes the program as 80386
assembly language for the InitechOS run-time library.  The language it
accepts is a compact Pascal designed to be exactly as large as a compiler
needs in order to compile itself: Turbo Initech is written in the language
this book describes, and it compiles its own source.

The Turbo Initech language provides:
\begin{itemize}
\item 32-bit signed \rw{integer} values, \rw{boolean} and \rw{char} values,
  and strings of up to 255 characters;
\item constants, static arrays, and record types with integer, boolean and
  character fields;
\item \rw{if}, \rw{while}, \rw{for}, and \rw{repeat} statements and
  compound statements;
\item procedures and functions with value and \rw{var} parameters, local
  variables, recursion, and \rw{forward} declarations;
\item complete evaluation of every Boolean expression;
\item untyped files read and written in blocks of bytes through InitechDOS.
\end{itemize}

Every program is one source file; all storage is static or is allocated in
the frame of the routine that declares it.  Turbo Initech has no pointers,
no sets, no enumerated types, no real numbers, no units and no \rw{case}
statement.  The word \rw{case} is reserved so that no program can use it as
an identifier.  These limits are part of the definition, and Chapter 3,
``Types,'' shows the constructions Turbo Initech programs use in their
place.

\section{The Turbo Initech manuals}
\ix{manuals}
Turbo Initech comes with four books and a Quick Reference.  Each book has
one job.

\SideHead{User's Guide}The \doctitle{Turbo Initech User's Guide}
(1077-0000-10) shows you how to run the compiler and teaches the language
by building programs, from the first one to programs that read and write
files.  It also describes the compiler's command and everything it prints.

\SideHead{Programmer's Guide}This book (1077-0000-11) defines the language,
the standard routines as a group, the internal representation of data, the
calling conventions, the code the compiler writes, and every compiler
message.

\SideHead{Library Reference}The \doctitle{Turbo Initech Library Reference}
(1077-0000-12) describes every standard procedure and function in
alphabetical order, each with an example program, and the entry points of
the run-time library.

\SideHead{Self-Compilation Guide}The \doctitle{Turbo Initech
Self-Compilation Guide} (1077-0000-13) describes how Turbo Initech is built
from its own source and the acceptance procedure that proves a new
compiler reproduces itself.

\SideHead{Quick Reference}The \doctitle{Turbo Initech Quick Reference}
(1077-0000-20) summarizes the syntax, the standard routines and the
messages.

Read the \doctitle{User's Guide} if you want to learn the language or to
know how to compile and run a program.  Read the \doctitle{Library
Reference} if you want the exact behavior of one standard routine.  Read the
\doctitle{Self-Compilation Guide} if you build Turbo Initech itself.

\section{Installing Turbo Initech}
\ix{installing}
The compiler is the single program \file{TPS.COM}.  Turbo Initech programs
are compiled with it and run under InitechDOS.  The \doctitle{User's Guide},
Chapter 1, describes the files you need and how to compile and run your
first program.
\draft{the Turbo Initech distribution diskette and its installation are not
produced (DS row TIN; PUB-REF-0001 Sec 12); TPS.COM is built by the
genesis compiler as the Self-Compilation Guide describes}

\section{Typefaces used in these books}
\ix{typefaces}
\ix{conventions}
The Turbo Initech manuals use these typefaces to mean these things:

\begin{BorTable}{Typefaces}{p{38mm}p{62mm}}
\Hd{Example} & \Hd{Meaning}\EndHd
\texttt{Writeln(Total);} & Program text, examples, and what the compiler and
  your programs print\\
\texttt{\bfseries begin} & A reserved word inside program text\\
\textbf{begin} & A reserved word in running text\\
\textit{Count} & An identifier of your program, a placeholder in a syntax
  description, or a standard routine named in running text\\
\textbf{\texttt{TPS}} & Something you type\\
\key{Enter} & A key on the keyboard\\
\MarginFace Margin note & A note that comments on the paragraph
  beside it
\end{BorTable}

Syntax is given by syntax diagrams; Chapter 1 explains how to read them.
Numbers in program text are decimal.

\section{How to contact Initech}
\ix{technical support}
If you have a question about Turbo Initech that this book does not answer,
call the Initech Help Desk at 555-0125, extension 2504, or write to
\InitechAddressLine.  Please have the part number of your manual and the
text of any message the compiler printed at hand.  Your Initech
representative can order further copies of the manuals for you.
